\chapter{Introduction}

\section{Présentation du projet}
\textbf{AcadémiePlus} est une plateforme web de soutien scolaire spécialisée en
\textbf{mathématiques}, destinée aux élèves du système éducatif algérien
(du primaire jusqu'au baccalauréat). L'application combine du contenu de cours
structuré, des exercices et quiz interactifs, un \textbf{moteur d'apprentissage
adaptatif} qui ajuste la difficulté au niveau réel de l'élève, et un
\textbf{assistant pédagogique IA} (chatbot) encadré par le programme scolaire.

\begin{infobox}[\faBullseye\ Objectif pédagogique]
Offrir à chaque élève un parcours \emph{personnalisé} : le système mesure en
permanence le niveau de l'élève (sur 100), détecte ses lacunes, génère des
exercices calibrés et fournit des commentaires et recommandations automatiques.
Les parents suivent la progression de leurs enfants via un tableau de bord dédié.
\end{infobox}

\section{Publics cibles et rôles}
La plateforme distingue quatre rôles applicatifs :
\begin{itemize}[leftmargin=1.4cm]
  \item[\faUserGraduate] \textbf{Élève (\texttt{student})} --- suit les cours, passe le test
        de positionnement, réalise exercices/quiz, dialogue avec l'IA.
  \item[\faUsers] \textbf{Parent (\texttt{parent})} --- relie ses enfants à son compte et
        consulte leurs rapports de progression.
  \item[\faChalkboardTeacher] \textbf{Pédagogue (\texttt{pedago})} --- rédige et gère le
        contenu des cours, exercices, quiz et examens (espace éditorial).
  \item[\faUserShield] \textbf{Administrateur (\texttt{admin})} --- supervise les
        utilisateurs, les abonnements, les codes d'activation et l'ensemble du contenu.
\end{itemize}

\section{Niveaux scolaires couverts}
Le contenu est organisé par niveau scolaire (énumération \texttt{school\_level}) :
\begin{center}
\begin{tabular}{ll}
\toprule
\textbf{Code interne} & \textbf{Niveau} \\
\midrule
\texttt{5eme\_primaire} & 5\textsuperscript{e} année primaire \\
\texttt{1ere\_cem} & 1\textsuperscript{re} année CEM (collège) \\
\texttt{2eme\_cem} & 2\textsuperscript{e} année CEM \\
\texttt{3eme\_cem} & 3\textsuperscript{e} année CEM \\
\texttt{4eme\_cem} & 4\textsuperscript{e} année CEM (BEM) \\
\texttt{premiere} & 1\textsuperscript{re} année secondaire \\
\texttt{seconde} & 2\textsuperscript{e} année secondaire \\
\texttt{terminale} & 3\textsuperscript{e} année secondaire (BAC) \\
\bottomrule
\end{tabular}
\end{center}

Chaque niveau secondaire peut être subdivisé en \textbf{filières} (sciences,
lettres, gestion, math-techniques\dots), ce qui détermine les chapitres affichés.

\section{Comment lire cette documentation}
\begin{itemize}
  \item Le \textbf{chapitre 2} décrit l'architecture globale.
  \item Le \textbf{chapitre 3} détaille la pile technique et les versions.
  \item Le \textbf{chapitre 4} explique l'authentification et la gestion des rôles.
  \item Le \textbf{chapitre 5} liste les fonctionnalités (pages) écran par écran.
  \item Le \textbf{chapitre 6} décrit le moteur adaptatif (calcul du niveau).
  \item Les \textbf{chapitres 7 et 8} documentent la base de données et le MCD.
  \item Le \textbf{chapitre 9} décrit les fonctions serveur (Edge Functions IA).
  \item Les \textbf{chapitres 10 et 11} traitent la sécurité et le déploiement.
\end{itemize}
